\documentclass[12pt]{article}
\usepackage[utf8]{inputenc}
\usepackage[T1]{fontenc}
\usepackage{hyperref}
\title{Tukudede Language Resources for K-12 Curriculum Development}
\author{Compiled by Claude Sonnet 5 (claude-sonnet-5) for Andr\'as Kornai}
\date{}
\begin{document}
\maketitle

\section{Introduction}
Tukudede (also Tokodede, Tocodede, Tukude; ISO 639-3 tkd) is an Austronesian (Timor group) language of Liqui\c{c}\'a municipality, Timor-Leste, around Maubara and Liqui\c{c}\'a. Wikipedia gives about 40,000 speakers (2010 census) and another source gives about 65,000. Both note decline. The search results were thin, so this dossier is short.

\section{Primary Educational Resources}
No primers, textbooks or school programmes were found.

\section{Linguistic and Reference Materials}
\begin{itemize}
\item Wikipedia, ``Tokodede language'': \url{https://en.wikipedia.org/wiki/Tokodede_language}. Its external links reportedly include an online Tokodede dictionary and a 2005 literary text. I did not see those pages themselves and cannot vouch for their scope.
\item Lexvo entry for tkd: \url{https://finto.fi/lexvo/en/page/tkd}.
\end{itemize}

\section{Religious and Community Resources}
\begin{itemize}
\item ScriptureEarth page for Tukudede: \url{https://www.scriptureearth.org/00eng.php?iso=tkd}. It lists a Tukudede Bible, with a pointer to Bible.com and audio from Global Recordings Network.
\end{itemize}

\section{Digital and Audio Resources}
The audio is on the ScriptureEarth page above. No corpus or software was found.

\section{Curriculum Development Context}
No school use was found. Portuguese and Tetun dominate formal education in Timor-Leste. A mother-tongue curriculum would have to build on the Bible translation, the online dictionary (to be verified) and community documentation.
\end{document}